\documentclass[12pt]{article}

%% Basic document formatting
\usepackage{amsmath, amsthm, amssymb, amsfonts}
\usepackage{mathtools}
\usepackage{mathrsfs}
\usepackage{xspace}
\usepackage{thmtools}
\usepackage{graphicx}
\usepackage{tikz-cd}
\usepackage{color}
\usepackage{setspace}
\usepackage{fancyhdr}
\usepackage{titling}
\usepackage{hyperref}
\usepackage[left=0.4in,right=0.4in,top=1in,bottom=1in]{geometry}
\usepackage{float}
\usepackage{tabularx}
\usepackage{booktabs}
\usepackage[utf8]{inputenc}
\usepackage[english]{babel}
\usepackage{framed}
\usepackage[dvipsnames]{xcolor}
\usepackage{environ}
\usepackage{tcolorbox}
\tcbuselibrary{theorems,skins,breakable}

\usepackage{enumitem}
\setlist[enumerate]{leftmargin=*}
\setlist[enumerate,1]{labelindent=\parindent}
\setlist[enumerate,2]{labelindent=0pt}

% Blackboard Bold
\newcommand{\Z}{\mathbb{Z}}                 % integers
\newcommand{\Zpl}{\mathbb{Z}_{+}}           % positive integers
\newcommand{\Znn}{\mathbb{Z}_{\geq 0}}      % nonnegative integers. 
\newcommand{\Q}{\mathbb{Q}}                 % rationals
\newcommand{\Qpl}{\mathbb{Q}_{+}}           % positive Rationals
\newcommand{\R}{\mathbb{R}}                 % reals
\newcommand{\Rpl}{\mathbb{R}_{+}}           % positive Reals
\newcommand{\C}{\mathbb{C}}                 % complex numbers
\newcommand{\F}{\mathbb{F}}                 % field

% Words
\newcommand{\st}{\text{ such that }}
\newcommand{\wrt}{\text{ with respect to }}
\newcommand{\with}{\text{ with }}
\newcommand{\ie}{\text{, i.e. }}
\newcommand*{\ditto}{\text{---}\texttt{"}\text{---}}

% Operators
\newcommand{\abs}[1]{\left|#1\right|}                   % absolute value
\newcommand{\norm}[1]{\abs{\abs{#1}}}
\newcommand{\floor}[1]{\left\lfloor #1 \right\rfloor}   % floor
\newcommand{\ceil}[1]{\left\lceil #1 \right\rceil}      % ceiling
\newcommand{\powset}[1]{\mathscr{P}(#1)}

% Category Theory 
\renewcommand{\hom}{\text{Hom}}
\newcommand{\homspace}[3]{\hom_{#1}(#2, #3)}
\newcommand{\coproduct}[9]{
  \begin{tikzcd}[ampersand replacement=\&]
      \& \& #1 \arrow[dll, bend right, "#6"] \arrow[dl, "#5"] \\
   #4 \& #3 \arrow[l, "#9"] \\
      \& \& #2 \arrow[ull, bend left, "#8"] \arrow[ul, "#7"]
  \end{tikzcd}
}

\newcommand{\product}[9]{ 
  \begin{tikzcd}[ampersand replacement=\&]
      \& \& #3 \\
      #1 \arrow[r, "#7"] \arrow[urr, bend left, "#5"] \arrow[drr, bend right, "#6"] \& #2 \arrow[ur, "#8"] \arrow[dr, "#9"] \\ 
      \& \& #4
  \end{tikzcd}
}


% Algebra 
\newcommand{\Syl}[2]{\text{Syl}_{#1}{(#2)}}           % set of Sylow-p subgroups 
\newcommand{\<}{\langle}                % \<x\>, subgroup generated by x 
\renewcommand{\>}{\rangle}
\renewcommand{\index}[2]{[#1 : #2]}
\newcommand{\id}{\text{id}}             % identity element
\newcommand{\lhdneq}{%
  \mathrel{\ooalign{$\lneq$\cr\raise.22ex\hbox{$\lhd$}\cr}}}
\newcommand{\order}[1]{\text{o}(#1)}    % order of an element
\newcommand{\spec}{\operatorname{Spec}}
\newcommand{\rad}{\operatorname{rad}}   % radical of an ideal 
\newcommand{\sym}[1]{\mathfrak{S}_{#1}}
\newcommand{\aut}[1]{\text{Aut}(#1)} 
\newcommand{\gal}[2]{\text{Gal}(#1 / #2)} 
\newcommand{\inn}[1]{\text{Inn}(#1)} 
\let\oldcong\cong
\let\oldequiv\equiv
\renewcommand{\cong}{\oldequiv}
\renewcommand{\equiv}{\oldcong}

\newcommand{\triangleleftneq}{\mathrel{\ooalign{%
  $\trianglelefteq$\cr
  \hidewidth\raise0.06ex\hbox{$\not\mkern4mu$}\hidewidth\cr
}}}

\makeatletter % cycle
\newcommand{\cyc}[1]{(\mathbf{\cyc@process#1\relax})}
\def\cyc@process#1#2\relax{%
	#1%
	\ifx\relax#2\relax
	\else
		\,\cyc@process#2\relax
	\fi
}
\makeatother

\newcommand{\GL}[2]{\text{GL}_{#1}(#2)}                             % general linear group
\newcommand{\SL}[2]{\text{SL}_{#1}(#2)}                             % special linear group
\newcommand{\gl}[2]{\mathfrak{gl}_{#1}(#2)}
\renewcommand{\sl}[2]{\mathfrak{sl}_{#1}(#2)}
\newcommand{\so}[2]{\mathfrak{so}_{#1}(#2)}
\newcommand{\su}[1]{\mathfrak{su}(#1)}

% Linear Algebra
\newcommand{\bmat}[1]{\begin{bmatrix}#1\end{bmatrix}}   % bracketed matrix
\newcommand{\rank}{\operatorname{rank}}                 % rank
\newcommand{\nullity}{\operatorname{nullity}}           % nullity
\newcommand{\tr}{\operatorname{tr}}

% Combinatorics
\newcommand{\qnum}[1]{\left[#1\right]_q}                % q-analog
\newcommand{\qfact}[1]{\left[#1\right]!_q}              % q-analog factorial 
\newcommand{\qbinom}[2]{\binom{#1}{#2}_q}               % Gaussian binomial coefficient

% Topology/Analysis
\newcommand{\ball}[2]{\text{B}_{#1}(#2)}  % B_r(x): open r-balls around x
\newcommand{\diam}{\text{diam}}           % diamter of a set in metric space 
\newcommand{\lowint}[2]{\underline{\int_{#1}^{#2}}}
\newcommand{\upint}[2]{\overline{\int}_{#1}^{#2}}

% Misc Notation
\newcommand{\oneton}[1]{\{1, \ldots, #1\}}
\newcommand{\defeq}{\vcentcolon=}   % :=
\newcommand{\eqdef}{=\vcentcolon}   % =: 
\renewcommand{\bf}[1]{\textbf{#1}}
\newcommand{\im}{\operatorname{im}}
\newcommand{\fnrest}[2]{\left.#1\right|_{#2}}
\newcommand{\isom}{\overset{\sim}{\rightarrow}} 


\renewcommand{\thesection}{\arabic{section}.}
\renewcommand{\thesubsection}{(\alph{subsection})}

\newtheorem{claim}{Claim}
\newtheorem*{lemma}{Lemma}

%% Headers & title setup
\newcommand{\course}{Math100C}
\newcommand{\myname}{Jay Ser}
\setlength{\headheight}{14.5pt}
\pagestyle{fancy}
\fancyhf{}
\renewcommand{\headrulewidth}{0.4pt}
\lhead{\course}
\rhead{\myname}
\cfoot{\thepage}
\setlength{\droptitle}{-4em} 
\title{\course\ - WK10} 
\author{\myname}
\date{2026.06.09}

\begin{document}
\maketitle
\thispagestyle{fancy}


%------------------------------------------------------------------------------%

\section{} 
Let $G$ be a group of order $p^n$ for a prime $p$ and integer $n > 1$. 
Prove $G$ has a normal subgroup of order $p$. 

\noindent \dotfill 

Consider the class equation for $G$: 
$$\#G = \#Z(G) + \sum_i \# \mathcal{O}(x_i)$$
The size of the nontrivial orbits (i.e., orbits of elements not in the center) each divide $\#G = p^n$, i.e., 
$p \mid \#\mathcal{O}(x_i)$ for each $i$. 
Suppose $G$ has a trivial center. 
Then 
$$p^n - 1 = \sum_{i} \# \mathcal{O}(x_i)$$
and $p$ divides the right hand side by the observation above. 
But clearly $p$ does not divide the left hand side, a contradition. 
It must be that $G$ has a nontrivial center; 
by the order of the group, $p \mid \#Z(G)$. 
Namely, there is an element of order $p$ in the center; 
this element generates a normal subgroup of $G$ of order $p$. 

\section{} 
Let $K / F$ be a Galois extension and $p$ any prime number. 
Show there exists an intermediate field $K / E / F$ such that 
\begin{itemize}
  \item $E / F$ has degree not divisible by $p$. 
  \item $K / E$ is Galois with degree a power of $p$. 
\end{itemize}

\noindent \dotfill 

Given a Galois extension $K / F$, recall that for any intermediate field $K / E / F$, $K / E$ is Galois. 
If $p \nmid \index{K}{F}$, then just take $E \defeq K$. 
Assume $\index{K}{F} = p^n a$ for $p \nmid a$ and $n \geq 1$. 
By Sylow's first theorem, there is a subgroup $H < \gal{K}{F}$ of order $p^n$. 
Using the Galois correspondence, one immediately checks that the fixed field of $H$ satisfies the desired properties. 


\section{} 
Show that every finite extension of $\R$ other than itself has even degree. 
Then, together with the previous problems, prove that every finite extension of $\C$ is trivial; 
in other words, every polynomial over $\C$ splits completely.  

\noindent \dotfill 

Take any finite, nontrivial extension $K / \R$. 
By the primitive element theorem, there is an element $\gamma \in K$ such that $K = \R(\gamma)$. 
$\R \subsetneq K$ shows $\gamma$ is of degree greater than 1 over $\R$. 
Since every non-linear odd-degree polynomimal over $\R$ has a root, the minimal polynomial for $\gamma$ must be of even degree. 
Hence $\index{K}{\R}$ is even. 

Suppose there exists a finite, nontrivial extension $K / \C$. 
Since $\C = \R(i)$ is a degree two extension of $\R$, $K / \R$ is a finite extension. 
Let $L / K$ be the Galois closure of $K$ with respect to $\R$. 
Applying the case $p = 2$ to the result of Q2, there is an intermediate field $L / E$ containing $\R$ such that $\index{L}{E} = 2^{t}$ and $2 \nmid \index{E}{\R}$. 
As shown above, the only finite extension of $\R$ of degree is the trivial extension; 
it follows $\index{L}{\R} = 2^t$, where $t > 1$ by assumption on $L, K$ and $\C$. 

If $t = 2$, i.e., $\index{L}{\R} = 4$, then namely $L$ is a degree 2 extension of $\C$, i.e., 
$L = \C(\sqrt{a})$ for some $a \in \C$. 
But De Moivre's theorem says $\sqrt{a} \in \C$, hence $L = \C$, a contradiction. 
Hence assume $t > 2$.  
By Q1, the Galois group $\gal{L}{\R}$ of order $2^{t}$ has a normal subgroup of order 2, thus yielding a degree $2^{t - 1}$ Galois extension of $\R$. 
In fact, repeating this argument gives Galois extensions $L_0 = L / L_1 / L_2 / \ldots / L_{t} = \R$ of degrees $2^t, 2^{t - 1}, 2^{t - 2}, \ldots, 2^{0} = 1$, respectively. 
Now, the $L_{t - 1}$, the degree 2 extension of $\R$, is just $\C$ since every degree two extension of $\R$ is $\C$; 
$\C$ itself is degree 2 over $\R$, and if $\R(\sqrt{-a})$, $a \in \R_{+}$, is any degree two extension of $\R$, then $\R(\sqrt{-a})$ contains $\sqrt{-a} \sqrt{1 / a} = i$ and therefore $\C$ as well. 
But then $L_{t - 2}$ is a degree 2 extension of $\C$, once again contradicting the fact that $\C$ doesn't have any quadratic extensions. 

The conclusion is that there are no nontrivial finite extensions $K / \C$. 
Namely, this shows every polynomial over $\C$ splits: 
if an irreducible polynomial over $\C$ has degree greater than 1, then its furnishes a nontrivial finite extension of $\C$, a contradiction. 

\section{} 
Prove every abelian group occurs as the Galois group of some Abelian extension of $\Q$. 

\noindent \dotfill 

I follow a line of reasoning sketched out in Dummit and Foote. 
Let $G$ be a finite abelian group. 
By the classification theorem for finitely generated modules over PIDs, namely the PID $\Z$, 
$$G \equiv (\Z / n_1 \Z) \times (\Z / n_2 \Z) \times \cdots \times (\Z / n_k \Z)$$
I will just write $Z_{n_i}$ in place of $\Z / n_i \Z$. 
By the (true) assumption that there are infinitely many primes congruent to 1 mod any integer $m$, one can take distinct primes $p_1, \ldots, p_k$ such that 
\begin{gather*}
  p_1 \cong 1 \mod n_1 \\ 
  \vdots  \\ 
  p_k \cong 1 \mod n_k
\end{gather*}
Since the primes $p_1, \ldots, p_k$ are distinct (namely, $p_i, p_j$ are coprime for $i \neq j$), one can apply the Chinese Remainder Theorem: 
letting $n = p_1 \cdots p_k$, 
$$\Z / n\Z \equiv (\Z / p_1 \Z) \times \cdots \times (\Z / p_k \Z)$$
as rings. 
Of course, this isomorphism of rings induces an isomorphism between the group of units: 
\begin{align*}
  (\Z / n\Z)^\times & \equiv (\Z / p_1 \Z)^{\times} \times \cdots \times (\Z / p_k \Z)^{\times} \\ 
					& \equiv Z_{p_1 - 1} \times \cdots \times Z_{p_k - 1}
\end{align*}
where the last isomorphism recalls the fact that the multiplicative group of $\Z / pZ$ for any prime $p$ is cyclic. 

Now, by construction, for each $i \in \oneton{k}$, $n_i \mid p_i - 1$. 
Recall that abelian groups have a subgroup of any order that divides the order of the group. 
Hence let $H_i < Z_{p_i - 1}$ be a subgroup of order $\frac{p_i - 1}{n_i}$; 
subgroups and quotients of cyclic groups are cyclic, so $Z_{p_i - 1} / H_i \equiv Z_{n_i}$. 
This identification gives the isomorphism 
\begin{align*}
  (\Z / n\Z)^\times / (H_1 \times \cdots \times H_k) 
  & \equiv (Z_{p_1 - 1} / H_1) \times \cdots \times (Z_{p_k - 1} / H_k) \\ 
  & \equiv Z_{n_1} \times \cdots \times Z_{n_k}  \\ 
  & \equiv G
\end{align*}

Finally, to realize $G$ as the Galois group of an abelian extension over $\Q$, recall, for any $m$, $\Q(\zeta_m) / \Q$, where $\zeta_m = e^{2\pi i / m}$ denotes the primitive $m$th root of unity, is a Galois extension with Galois group $(\Z / m\Z)^{\times}$. 
Also recall that if $a, b$ are coprime, then $\Q(\zeta_a, \zeta_b) \equiv \Q(\zeta_{ab})$. 
From these facts, one can see 
$$\Q(\zeta_{p_1}, \ldots, \zeta_{p_k}) \equiv \Q(\zeta_n)$$
and this field has Galois group $(\Z / n\Z)^\times$ over $\Q$. 
Furthermore, the Galois group has a (normal) subgroup isomorphic to $H \defeq H_1 \times \cdots \times H_k$, and thus the Galois correspondence says the fixed field $\Q(\zeta_n)^{H}$ is an (abelian) Galois extension of $\Q$ with Galois group $G$. 

\end{document}
